\chapter{Prompts del modelo de lenguaje}
\label{anexo:prompts}

Este anexo documenta los dos usos del modelo Llama~3.3-70b en el
framework: (i) la curación clínica de metadatos de GEO, y (ii) el
asistente conversacional integrado en la interfaz. Ambos usos son bien
delimitados y prohíben explícitamente al modelo desbordar su función.

\section*{Prompt de curación clínica}
\phantomsection
\addcontentsline{toc}{section}{Prompt de curación clínica}

El módulo \texttt{paso2\_metadata.py} construye para cada muestra un
prompt del tipo:

\begin{lstlisting}[caption={Prompt de curación (esquemático).}]
Eres un curador clinico automatizado. Tu tarea es leer los metadatos
crudos de una muestra depositada en NCBI GEO y clasificarla en UNO de
los tres grupos siguientes:

  1. sano       -> tejido pulmonar no tumoral (control adyacente o
                   pulmon sano)
  2. enfermo    -> tejido tumoral primario de pulmon (adenocarcinoma,
                   carcinoma escamoso, etc.)
  3. ambiguo    -> caso que no permite decidir sin equivocarse (linea
                   celular sin condicion, histologia rara, ausencia de
                   informacion, etc.)

Reglas:
- Responde SOLO con JSON valido de la forma
  {"etiqueta": "sano"|"enfermo"|"ambiguo", "motivo": "..."}
- El campo "motivo" debe ser una frase corta que cite las palabras del
  texto crudo que motivan la clasificacion.
- Prefiere "ambiguo" antes que forzar una etiqueta incierta.
- No inventes informacion que no este en el texto.

METADATOS CRUDOS DE LA MUESTRA:
  title:                {title}
  source_name:          {source_name}
  characteristics_ch1:  {characteristics_ch1}
\end{lstlisting}

El prompt se instancia con los campos de cada muestra rellenados. Se
solicita al modelo con \texttt{temperature=0.0} para maximizar el
determinismo. La respuesta se parsea como JSON y se rechaza (dejando
etiqueta \texttt{ambiguo}) si el parseo falla.

La tasa de éxito global de la curación es del 85,9\,\% de media sobre
las 11 cohortes descargadas, con 9 cohortes por encima del 95\,\%.
GSE40791 es la única cohorte que queda por debajo del umbral de
inclusión (curación por debajo del 20\,\%), debido a la excepcional
falta de detalle en sus metadatos originales.

\section*{Prompt sistema del asistente conversacional}
\phantomsection
\addcontentsline{toc}{section}{Prompt sistema del asistente conversacional}

El asistente conversacional integrado en la interfaz emplea un prompt
sistema construido dinámicamente por \texttt{contexto\_tfm.py}. La
estructura consta de tres bloques: (i) declaración de rol y ámbito, (ii)
reglas de comportamiento en orden de prioridad, y (iii) el contexto de
resultados leído en tiempo real de los CSV/JSON del pipeline. La
figura~\ref{fig:prompt} representa gráficamente cómo se compone el
prompt en cada ejecución del asistente.

\begin{figure}[H]
  \centering
  \begin{tikzpicture}[node distance=0.55cm and 0.9cm]
    % Fuentes de datos (izquierda)
    \node[entrada, minimum width=3.4cm, minimum height=0.75cm] (aud)   {\footnotesize AUDITORIA\_COHORTES.csv};
    \node[entrada, below=of aud, minimum width=3.4cm, minimum height=0.75cm] (lodo)  {\footnotesize LODO\_HONESTO\_RESULTADOS.csv};
    \node[entrada, below=of lodo, minimum width=3.4cm, minimum height=0.75cm] (sub)   {\footnotesize SUBTIPO\_LODO\_RESULTADOS.csv};
    \node[entrada, below=of sub, minimum width=3.4cm, minimum height=0.75cm] (comp)  {\footnotesize COMPOSICION\_VS\_BIOLOGIA.csv};
    \node[entrada, below=of comp, minimum width=3.4cm, minimum height=0.75cm] (rf)    {\footnotesize FIRMA\_VALIDADA\_RESUMEN.json};
    \node[entrada, below=of rf, minimum width=3.4cm, minimum height=0.75cm] (comp2) {\footnotesize FIRMA\_VALIDADA\_COMPLETA.csv};
    % Modulo
    \node[paso, right=of lodo, yshift=-1.1cm, minimum width=3.6cm, minimum height=1.6cm] (mod) {\texttt{contexto\_tfm.py}\\\footnotesize construir\_contexto()\\\footnotesize \_cifras\_clave()};
    % Prompt bloques
    \node[salida, right=of mod, yshift=1.6cm, minimum width=4cm] (rol)    {\textbf{Rol y ámbito}\\\footnotesize (estático)};
    \node[salida, right=of mod, minimum width=4cm]              (reglas) {\textbf{8 reglas}\\\footnotesize (estático + cifras clave)};
    \node[salida, right=of mod, yshift=-1.6cm, minimum width=4cm] (ctx)  {\textbf{Contexto de resultados}\\\footnotesize (dinámico desde CSVs)};
    % LLM
    \node[paso, right=of reglas, minimum width=2.6cm, fill=purple!8, draw=purple!60!black] (llm) {\textbf{Llama 3.3-70b}\\\footnotesize vía Groq};
    % Historial
    \node[paso, above=of llm, minimum width=2.6cm] (hist) {Últimos 8\\mensajes};
    % Flechas
    \draw[flecha] (aud.east)  -- (mod.west);
    \draw[flecha] (lodo.east) -- (mod.west);
    \draw[flecha] (sub.east)  -- (mod.west);
    \draw[flecha] (comp.east) -- (mod.west);
    \draw[flecha] (rf.east)   -- (mod.west);
    \draw[flecha] (comp2.east) -- (mod.west);
    \draw[flecha] (mod.east) -- (rol.west);
    \draw[flecha] (mod.east) -- (reglas.west);
    \draw[flecha] (mod.east) -- (ctx.west);
    \draw[flecha] (rol.east) -- (llm.west);
    \draw[flecha] (reglas.east) -- (llm.west);
    \draw[flecha] (ctx.east) -- (llm.west);
    \draw[flecha] (hist) -- (llm);
  \end{tikzpicture}
  \caption{Composición del prompt sistema del asistente conversacional.
  Los CSV/JSON del pipeline (bloques verdes) alimentan
  \texttt{contexto\_tfm.py}, que compone tres secciones (rol, reglas y
  contexto dinámico) que se concatenan con los últimos ocho mensajes
  del historial y se envían al modelo Llama~3.3-70b vía Groq. Como los
  CSV/JSON se leen en tiempo real, cualquier reejecución del pipeline
  se propaga al chat sin editar código.}
  \label{fig:prompt}
\end{figure}

\subsection*{Reglas del sistema}
\phantomsection
\addcontentsline{toc}{subsection}{Reglas del sistema}

Las ocho reglas que gobiernan el comportamiento del asistente son:

\begin{enumerate}
  \item \textbf{Fidelidad a los datos.} Responder sólo con información
  del contexto. Si algo no está, decir literalmente ``Eso no figura en
  los resultados del trabajo''. No completar con conocimiento general ni
  estimar cifras.

  \item \textbf{Enfoque.} El trabajo es de análisis de datos y ML
  aplicado a transcriptómica. Priorizar los resultados biológicos y de
  rendimiento del modelo, no la metodología de calidad de datos.

  \item \textbf{Resultados a citar cuando encajen.} Firma de 1174 genes,
  panel mínimo de 20 genes con AUC 0,966 vs 0,970 de la firma completa,
  18/20 marcadores IHC recuperados (12/12 escamosos, 6/8 adenocarcinoma),
  AUC LODO tumor vs sano 0,925, AUC LODO subtipo 0,968. Top de la firma
  por linaje. \emph{No citar SFTPC ni SFTPA1 como parte de la firma}:
  son marcadores IHC pero no se recuperan.

  \item \textbf{Biología como fin.} Al preguntar por genes o por qué el
  modelo funciona, dar la interpretación biológica (desmosomas
  escamosos, queratinas basales, TP63 como factor maestro, NAPSA como
  aspartil-proteasa alveolar, SFTPB como surfactante, NKX2-1 como
  factor de transcripción adenocarcinoma).

  \item \textbf{Criterios de inclusión, no ``auditoría''.} Al preguntar
  por integridad de datos, enmarcarlo como criterios de inclusión del
  pipeline: 8 de 11 cohortes pasan los criterios, el filtrado es
  transparente y reproducible.

  \item \textbf{Métricas del ML.} Al preguntar por rendimiento, dar AUC,
  balanced accuracy, sensibilidad, especificidad. Aclarar que el umbral
  de decisión requiere calibrarse entre cohortes por desbalance de
  entrenamiento.

  \item \textbf{Sin consejo médico.} Aclarar que la firma no está
  validada para uso clínico.

  \item \textbf{Fuera de ámbito.} Si la pregunta no guarda relación con
  el trabajo ni con transcriptómica de cáncer de pulmón, responder
  ``Fuera de ámbito: esta consulta no corresponde a los resultados de
  este TFM''.
\end{enumerate}

\subsection*{Contexto dinámico}
\phantomsection
\addcontentsline{toc}{subsection}{Contexto dinámico}

El bloque de contexto se genera con \texttt{construir\_contexto()} y
contiene únicamente cifras leídas de los CSV/JSON del pipeline:
descripción del sistema, número de cohortes incluidas y muestras
evaluables, firma génica con desglose IHC, rendimiento ML de ambas
tareas, top-10 genes, panel de proliferación y panel de pulmón sano.
Estas cifras se recomputan en cada arranque del asistente, por lo que
cualquier reejecución del pipeline se propaga automáticamente al chat.

En particular, las listas de marcadores IHC recuperados y no recuperados
se calculan por \emph{intersección real} entre las constantes
\texttt{IHC\_CLINICA} (anexo~\ref{anexo:paneles}) y el CSV
\texttt{FIRMA\_VALIDADA\_COMPLETA.csv}. Esto elimina la posibilidad de
que un cambio en la firma deje al asistente afirmando cosas incorrectas
sobre marcadores.

\subsection*{Ejemplos de interacción}
\phantomsection
\addcontentsline{toc}{subsection}{Ejemplos de interacción}

Algunos ejemplos representativos de respuestas del asistente:

\begin{description}
  \item[Usuario:] ¿De qué trata el proyecto?
  \item[Asistente:] Es un framework de análisis transcriptómico y ML
  para identificar biomarcadores de cáncer de pulmón. Integra cohortes
  públicas de NCBI GEO, cura los metadatos con Llama 3.3-70b, entrena
  clasificadores supervisados sobre dos tareas (tumor vs sano y subtipo
  histológico) y entrega, para el subtipo, una firma de 1174 genes
  replicados y un panel mínimo de 20 genes con AUC 0,966.
  18/20 marcadores IHC diagnósticos se recuperan de novo.

  \item[Usuario:] ¿Qué es SFTPC?
  \item[Asistente:] SFTPC es una proteína del surfactante alveolar
  (surfactant protein C), considerada marcador canónico del programa
  alveolar tipo II. Forma parte del panel IHC clínico de adenocarcinoma
  (OMS 2015), pero \emph{no está en la firma génica validada} del
  framework: es uno de los dos marcadores del panel ADC (junto con
  SFTPA1) que no supera el criterio de replicación en las tres cohortes
  independientes.

  \item[Usuario:] ¿Cuál es el AUC del clasificador tumor vs sano?
  \item[Asistente:] La AUC media LODO sobre las 8 cohortes evaluables
  es 0,925, con balanced accuracy 0,772, sensibilidad 0,983 y
  especificidad 0,561. La brecha entre AUC y balanced accuracy proviene
  del desbalance del entrenamiento (938 tumores frente a 219 controles);
  se corrige recalibrando el umbral por cohorte.
\end{description}
